% GENERATED FILE. Edit canonical lessons, then run npm run generate:books.
% canonical-source-hash: fnv1a64:3fa536a348754fd7
% canonical-lessons: PA-C234-nind, PA-C234-diva, PA-C234-balna, PA-C234-bujhauna, PA-C234-palangh

\chapter{Lights Out}
\label{ch:pa-a2-lights-out}
\hlchaptermodality{234}

\begin{chapteropening}
\textbf{By the end of this chapter:} \emph{I can say what I do before I sleep.}

Complete the last lesson of chapter 234: I can say what I do before I sleep.
\end{chapteropening}

\section[\texorpdfstring{\pa{ਨੀਂਦ} (nīnd)}{ਨੀਂਦ (nīnd)}]{\pa{ਨੀਂਦ} (nīnd) — sleep}

\label{lesson:PA-C234-nind}



\begin{quote}
Before the new one: say the Punjabi for a thick cotton sheet, then the Punjabi for a fan.
\end{quote}

\subsection*{You'll want to know: \pa{ਨੀਂਦ}}
\textbf{\pa{ਨੀਂਦ}} — \emph{nīnd} — \textquotedblleft{}sleep\textquotedblright{}.

\begin{quote}
\pa{ਰਾਤ} \pa{ਨੂੰ} \pa{ਮੈਨੂੰ} \pa{ਚੰਗੀ} \pa{ਨੀਂਦ} \pa{ਆਉਂਦੀ} \pa{ਹੈ।}
\end{quote}

\begin{quote}
\emph{rāt nū̃ mainū̃ chaṅgī nīnd āundī hai.}
\end{quote}

\begin{quote}
I sleep well at night.
\end{quote}

A feminine noun.

\begin{cousinweb}
\textbf{\pa{ਨੀਂਦ}} (\emph{nīnd}) is inherited from Sanskrit \emph{nidrā}, \textquotedblleft{}sleep\textquotedblright{}; CDIAL lists Punjabi \emph{nī̃d}, \textquotedblleft{}sleep\textquotedblright{} (CDIAL 7200).
\end{cousinweb}

\begin{grammarlens}[title={What you've built}]
One more word for saying what you do.
\end{grammarlens}

\subsection*{Your turn}
\begin{itemize}
\raggedright
  \item \emph{Say it:} \emph{nīnd}
  \item \emph{Say it:} \emph{nīnd}, once more
  \item \emph{Say it:} \emph{rāt nū̃ mainū̃ chaṅgī nīnd āundī hai.} — I sleep well at night.
  \item \emph{Recall:} say the Punjabi for a thick cotton sheet, then the Punjabi for a fan, then say \emph{nīnd} again
\end{itemize}

\begin{tcolorbox}[breakable,colback=teal!4,colframe=teal!35!black,title={Before you move on}]
What does \pa{ਨੀਂਦ} mean? (\textquotedblleft{}Sleep\textquotedblright{}.) Say it once more.
\end{tcolorbox}
\section[\texorpdfstring{\pa{ਦੀਵਾ} (dīvā)}{ਦੀਵਾ (dīvā)}]{\pa{ਦੀਵਾ} (dīvā) — an oil lamp}

\label{lesson:PA-C234-diva}



\begin{quote}
Before the new one: say the Punjabi for a fan, then the Punjabi for sleep (noun).
\end{quote}

\subsection*{You'll want to know: \pa{ਦੀਵਾ}}
\textbf{\pa{ਦੀਵਾ}} — \emph{dīvā} — \textquotedblleft{}an oil lamp\textquotedblright{}.

\begin{quote}
\pa{ਕਮਰੇ} \pa{ਵਿੱਚ} \pa{ਇੱਕ} \pa{ਛੋਟਾ} \pa{ਦੀਵਾ} \pa{ਹੈ।}
\end{quote}

\begin{quote}
\emph{kamre vich ikk chhoṭā dīvā hai.}
\end{quote}

\begin{quote}
There is a small oil lamp in the room.
\end{quote}

A masculine noun.

\begin{cousinweb}
\textbf{\pa{ਦੀਵਾ}} (\emph{dīvā}) is inherited from Sanskrit \emph{dīpa}, \textquotedblleft{}lamp\textquotedblright{}; CDIAL lists Punjabi \emph{dīvā}, \textquotedblleft{}lamp\textquotedblright{} (CDIAL 6348).
\end{cousinweb}

\begin{grammarlens}[title={What you've built}]
One more word for saying what you do.
\end{grammarlens}

\subsection*{Your turn}
\begin{itemize}
\raggedright
  \item \emph{Say it:} \emph{dīvā}
  \item \emph{Say it:} \emph{dīvā}, once more
  \item \emph{Say it:} \emph{kamre vich ikk chhoṭā dīvā hai.} — There is a small oil lamp in the room.
  \item \emph{Recall:} say the Punjabi for a fan, then the Punjabi for sleep (noun), then say \emph{dīvā} again
  \item \emph{Write it:} copy \textbf{\pa{ਦੀਵਾ}} once, checking the whole word
\end{itemize}

\begin{tcolorbox}[breakable,colback=teal!4,colframe=teal!35!black,title={Before you move on}]
What does \pa{ਦੀਵਾ} mean? (\textquotedblleft{}An oil lamp\textquotedblright{}.) Say it once more.
\end{tcolorbox}
\section[\texorpdfstring{\pa{ਬਲਣਾ} (balṇā)}{ਬਲਣਾ (balṇā)}]{\pa{ਬਲਣਾ} (balṇā) — to burn (of a flame)}

\label{lesson:PA-C234-balna}



\begin{quote}
Before the new one: say the Punjabi for sleep (noun), then the Punjabi for an oil lamp.
\end{quote}

\subsection*{You'll want to know: \pa{ਬਲਣਾ}}
\textbf{\pa{ਬਲਣਾ}} — \emph{balṇā} — \textquotedblleft{}to burn (of a flame)\textquotedblright{}.

\begin{quote}
\pa{ਦੀਵਾ} \pa{ਰਾਤ} \pa{ਨੂੰ} \pa{ਬਲਦਾ} \pa{ਹੈ।}
\end{quote}

\begin{quote}
\emph{dīvā rāt nū̃ baldā hai.}
\end{quote}

\begin{quote}
The lamp burns at night.
\end{quote}

Past (it …): \pa{ਉਹ} \pa{ਦੀਵਾ} \pa{ਬਲਿਆ}, \emph{oh dīvā baliā}.

\begin{cousinweb}
\textbf{\pa{ਬਲਣਾ}} (\emph{balṇā}) is inherited from a reconstructed \emph{*dvalati}, \textquotedblleft{}burns\textquotedblright{}; CDIAL lists Punjabi \emph{balṇā}, \textquotedblleft{}to burn, blaze\textquotedblright{} (CDIAL 6654).
\end{cousinweb}

\begin{grammarlens}[title={What you've built}]
One more word for saying what you do.
\end{grammarlens}

\subsection*{Your turn}
\begin{itemize}
\raggedright
  \item \emph{Say it:} \emph{balṇā}
  \item \emph{Say it:} \emph{balṇā}, once more
  \item \emph{Say it:} \emph{dīvā rāt nū̃ baldā hai.} — The lamp burns at night.
  \item \emph{Recall:} say the Punjabi for sleep (noun), then the Punjabi for an oil lamp, then say \emph{balṇā} again
  \item \emph{Write it:} copy \textbf{\pa{ਬਲਣਾ}} once, checking the whole word
\end{itemize}

\begin{tcolorbox}[breakable,colback=teal!4,colframe=teal!35!black,title={Before you move on}]
What does \pa{ਬਲਣਾ} mean? (\textquotedblleft{}To burn (of a flame)\textquotedblright{}.) Say it once more.
\end{tcolorbox}
\section[\texorpdfstring{\pa{ਬੁਝਾਉਣਾ} (bujhāuṇā)}{ਬੁਝਾਉਣਾ (bujhāuṇā)}]{\pa{ਬੁਝਾਉਣਾ} (bujhāuṇā) — to put out (a light or a fire)}

\label{lesson:PA-C234-bujhauna}



\begin{quote}
Before the new one: say the Punjabi for an oil lamp, then the Punjabi for to burn.
\end{quote}

\subsection*{You'll want to know: \pa{ਬੁਝਾਉਣਾ}}
\textbf{\pa{ਬੁਝਾਉਣਾ}} — \emph{bujhāuṇā} — \textquotedblleft{}to put out (a light or a fire)\textquotedblright{}.

\begin{quote}
\pa{ਮੈਂ} \pa{ਰਾਤ} \pa{ਨੂੰ} \pa{ਬੱਤੀ} \pa{ਬੁਝਾਉਂਦੀ} \pa{ਹਾਂ।}
\end{quote}

\begin{quote}
\emph{maiṁ rāt nū̃ battī bujhāundī hāṁ.}
\end{quote}

\begin{quote}
At night I put out the light.
\end{quote}

Past (I …): \pa{ਮੈਂ} \pa{ਦੀਵਾ} \pa{ਬੁਝਾਇਆ}, \emph{maiṁ dīvā bujhāiā}.

\begin{cousinweb}
\textbf{\pa{ਬੁਝਾਉਣਾ}} (\emph{bujhāuṇā}) is inherited from a reconstructed \emph{*vijjhāyati}, \textquotedblleft{}is burnt out, is extinguished\textquotedblright{}; CDIAL lists Punjabi \emph{bujhāuṇā}, \textquotedblleft{}to extinguish\textquotedblright{} (CDIAL 11703).
\end{cousinweb}

\begin{grammarlens}[title={What you've built}]
One more word for saying what you do.
\end{grammarlens}

\subsection*{Your turn}
\begin{itemize}
\raggedright
  \item \emph{Say it:} \emph{bujhāuṇā}
  \item \emph{Say it:} \emph{bujhāuṇā}, once more
  \item \emph{Say it:} \emph{maiṁ rāt nū̃ battī bujhāundī hāṁ.} — At night I put out the light.
  \item \emph{Recall:} say the Punjabi for an oil lamp, then the Punjabi for to burn, then say \emph{bujhāuṇā} again
  \item \emph{Write it:} copy \textbf{\pa{ਬੁਝਾਉਣਾ}} once, checking the whole word
\end{itemize}

\begin{tcolorbox}[breakable,colback=teal!4,colframe=teal!35!black,title={Before you move on}]
What does \pa{ਬੁਝਾਉਣਾ} mean? (\textquotedblleft{}To put out (a light or a fire)\textquotedblright{}.) Say it once more.
\end{tcolorbox}
\section[\texorpdfstring{\pa{ਪਲੰਘ} (palaṅgh)}{ਪਲੰਘ (palaṅgh)}]{\pa{ਪਲੰਘ} (palaṅgh) — a wooden bed}

\label{lesson:PA-C234-palangh}



\begin{quote}
Before the new one: say the Punjabi for to burn, then the Punjabi for to put out.
\end{quote}

\subsection*{You'll want to know: \pa{ਪਲੰਘ}}
\textbf{\pa{ਪਲੰਘ}} — \emph{palaṅgh} — \textquotedblleft{}a wooden bed\textquotedblright{}.

\begin{quote}
\pa{ਮੇਰਾ} \pa{ਪਲੰਘ} \pa{ਖਿੜਕੀ} \pa{ਦੇ} \pa{ਨੇੜੇ} \pa{ਹੈ।}
\end{quote}

\begin{quote}
\emph{merā palaṅgh khiṛkī de neṛe hai.}
\end{quote}

\begin{quote}
My bed is near the window.
\end{quote}

A masculine noun.

\begin{cousinweb}
\textbf{\pa{ਪਲੰਘ}} (\emph{palaṅgh}) is inherited from Sanskrit \emph{palyaṅka}, \textquotedblleft{}bed\textquotedblright{}; CDIAL lists Punjabi \emph{palaṅgh}, \textquotedblleft{}bed\textquotedblright{} (CDIAL 7964).
\end{cousinweb}

\begin{grammarlens}[title={What you've built}]
One more word for saying what you do.
\end{grammarlens}

\subsection*{Your turn}
\begin{itemize}
\raggedright
  \item \emph{Say it:} \emph{palaṅgh}
  \item \emph{Say it:} \emph{palaṅgh}, once more
  \item \emph{Say it:} \emph{merā palaṅgh khiṛkī de neṛe hai.} — My bed is near the window.
  \item \emph{Recall:} say the Punjabi for to burn, then the Punjabi for to put out, then say \emph{palaṅgh} again
  \item \emph{Write it:} copy \textbf{\pa{ਪਲੰਘ}} once, checking the whole word
\end{itemize}

\begin{tcolorbox}[breakable,colback=teal!4,colframe=teal!35!black,title={Before you move on}]
What does \pa{ਪਲੰਘ} mean? (\textquotedblleft{}A wooden bed\textquotedblright{}.) Say it once more.
\end{tcolorbox}
